% ============================================================
%  SECCIÓN — DESCRIPCIÓN DE OPCIONES Y SOLUCIÓN ADOPTADA
% ============================================================
\section{Análisis de opciones y solución adoptada}
\label{sec:muro_opciones}

A partir de los parámetros geotécnicos del Sondeo~1 y las condiciones
de contorno del proyecto (desnivel $H = \alturaDesnivelMuro$~m, longitud
150~ml, N.F.\ a $-2{,}50$~m, $3\,500$~mg~SO$_4^{2-}$/kg, zona
medioambientalmente protegida), se estudian las dos tipologías estructurales
posibles para la contención.

% -----------------------------------------------------------
\subsection{Opción A — Muro de contención en ménsula (hormigón armado)}
\label{subsec:opcion_mensula}
% -----------------------------------------------------------

\subsubsection{Descripción de la tipología}

El \textbf{muro de contención en ménsula} es una estructura de hormigón
armado compuesta por un \emph{fuste} en voladizo y una \emph{zapata corrida}
enterrada. El fuste actúa como ménsula empotrada en la zapata; el peso propio
del muro y el de las tierras situadas sobre el talón de la zapata son los
principales mecanismos estabilizadores. Se adopta la variante de \textbf{muro
en T invertida} (con talón largo hacia el trasdós), que optimiza la
utilización del peso del relleno para garantizar la estabilidad frente al
vuelco y al deslizamiento.

\subsubsection{Predimensionamiento geométrico — primer tanteo ($B = 4{,}00$~m)}
\label{subsubsec:predim_4m}

Se aplican las reglas empíricas habituales para muros en ménsula de hormigón
armado con altura $H = 7{,}30$~m. A continuación se justifica cada dimensión adoptada.

\paragraph{Espesores del fuste.}

\begin{align}
  e_1 &= \frac{H}{12} = \frac{7{,}3}{12} = 0{,}61~\text{m}
        \quad\longrightarrow\quad \resultado{e_1 = 0{,}60~\text{m}}
        \label{eq:e1_opc} \\[4pt]
  e_2 &= \frac{H}{10} = \frac{7{,}3}{10} = 0{,}73~\text{m}
        \quad\longrightarrow\quad \resultado{e_2 = 0{,}75~\text{m}}
        \label{eq:e2_opc}
\end{align}

El espesor en coronación ($e_1 = H/12$) es el mínimo necesario para alojar la
armadura con los recubrimientos reglamentarios (50~mm) y garantizar la
trabajabilidad del hormigón durante el vertido. El espesor en la base del fuste
($e_2 = H/10$) se incrementa para resistir el cortante y el momento flector
máximos, que aparecen en la sección de empotramiento con la zapata.

\paragraph{Canto de la zapata.}

\begin{equation}
  h_z = \frac{H}{10} \approx 0{,}73~\text{m}
        \quad\longrightarrow\quad \resultado{h_z = 0{,}80~\text{m}}
        \label{eq:hz_opc}
\end{equation}

Se redondea al alza hasta $h_z = 0{,}80$~m para cumplir la condición
$h_z \geq e_2$, que asegura el empotramiento correcto del fuste en la zapata
sin singularidades de armado ni concentraciones de tensión en la esquina.

\paragraph{Profundidad de cimentación.}

Se adopta $D_f = 1{,}00$~m medida desde la rasante del intradós hasta la
\textbf{base} de la zapata. Esta elección responde a tres criterios:
\begin{itemize}
  \item Supera holgadamente el mínimo reglamentario de $0{,}50$~m para
        cimentaciones superficiales.
  \item Protege la puntera de la zapata frente a posibles socavaciones
        superficiales y proporciona confinamiento suficiente para movilizar
        el empuje pasivo en el intradós.
\end{itemize}

\paragraph{Ancho total de la zapata — primer tanteo.}

La regla empírica propone:
\begin{equation}
  B \approx 0{,}55 \cdot H = 0{,}55 \times 7{,}30 = \resultado{B = 4{,}00~\text{m}}
  \label{eq:B_tanteo}
\end{equation}

Esta dimensión se adopta como \textbf{primer tanteo}, sujeta a verificación de
estabilidad al vuelco y al deslizamiento en los apartados siguientes.

\paragraph{Vuelo de puntera y ancho de talón.}

El vuelo de puntera se fija igual al espesor de la base del fuste
($B_\text{punta} \approx e_2$), valor mínimo que permite anclar correctamente
el armado inferior de la zapata:
\begin{equation*}
  B_\text{punta} = 0{,}80~\text{m}
  \qquad
  B_\text{talón} = B - e_2 - B_\text{punta}
    = 4{,}00 - 0{,}75 - 0{,}80 = \resultado{2{,}45~\text{m}}
\end{equation*}

Se adopta la variante en \textbf{T invertida} (talón largo hacia el trasdós):
el peso del relleno sobre el talón actúa como principal mecanismo estabilizador
frente al vuelco, minimizando la necesidad de aumentar el vuelo de puntera.

\begin{figure}[H]
  \centering
  \includegraphics[width=0.75\textwidth]{figures/muro/muro_z_4m}
  \caption{Geometría del muro en ménsula — primer tanteo ($B = 4{,}00$~m).}
  \label{fig:muro_z_4m}
\end{figure}

\begin{table}[H]
  \centering
  \caption{Resumen del predimensionamiento geométrico — primer tanteo}
  \label{tab:predim_muro}
  \resizebox{\textwidth}{!}{%
  \begin{tabular}{@{}llccp{6cm}@{}}
    \toprule
    \textbf{Dimensión} & \textbf{Símbolo}
      & \textbf{Regla} & \textbf{Valor} & \textbf{Criterio adoptado} \\
    \midrule
    Altura de contención
      & $H$ & — & 7{,}30~m
      & Desnivel impuesto por el proyecto (parámetro de grupo). \\[4pt]
    Espesor en coronación
      & $e_1$ & $H/12$ & 0{,}60~m
      & Mínimo constructivo: aloja armadura con recubrimiento de 50~mm. \\[4pt]
    Espesor en base del fuste
      & $e_2$ & $H/10$ & 0{,}75~m
      & Resistencia a cortante y flexión en la sección de empotramiento. \\[4pt]
    Canto de zapata
      & $h_z$ & $H/10$ & 0{,}80~m
      & Redondeo al alza para cumplir $h_z \geq e_2$ (empotramiento fuste). \\[4pt]
    Prof.\ cimentación (a base)
      & $D_f$ & — & 1{,}00~m
      & $> 0{,}50$~m reglamentario; deja 1{,}50~m sobre NF; protege la puntera. \\[4pt]
    Ancho total (tanteo)
      & $B$ & $0{,}55\,H$ & 4{,}00~m
      & Primer tanteo sujeto a verificación de estabilidad al vuelco. \\[4pt]
    Vuelo de puntera
      & $B_\text{punta}$ & $\approx e_2$ & 0{,}80~m
      & Mínimo para anclar el armado inferior de la zapata. \\[4pt]
    Ancho de talón
      & $B_\text{talón}$ & $B - e_2 - B_\text{punta}$ & 2{,}45~m
      & Talón largo (T invertida): estabilidad al vuelco por peso del relleno. \\
    \bottomrule
  \end{tabular}}
\end{table}

% -----------------------------------------------------------
\subsection{Fuerzas verticales y momentos estabilizadores — $B = 4{,}00$~m}
\label{subsec:fuerzas_verticales_4m}
% -----------------------------------------------------------

Se calculan los pesos por metro lineal de cada componente y su distancia $x$
a la punta de la zapata ($x = 0$ en el borde del intradós, positivo hacia el
trasdós). El brazo de los elementos situados sobre el talón vale:

\[
  x_\text{talón} = B_\text{punta} + e_2 + \frac{B_\text{talón}}{2}
    = 0{,}80 + 0{,}75 + \frac{2{,}45}{2} = \mathbf{2{,}78}~\text{m}
\]

\paragraph{A) Zapata ($B \times h_z$, $\gamma_\text{HA-30} = 2{,}50$~Tn/m³):}

\begin{align*}
  W_p &= \gamma_\text{HA} \cdot B \cdot h_z
       = 2{,}50 \times 4{,}00 \times 0{,}80
       = \resultado{8{,}00~\text{Tn/ml}} \\
  x_p &= \frac{B}{2} = \frac{4{,}00}{2} = 2{,}00~\text{m} \\
  M_p &= 8{,}00 \times 2{,}00 = \resultado{16{,}00~\text{Tn·m/ml}}
\end{align*}

\paragraph{B) Fuste (sección trapezoidal):}

Se descompone la sección en un rectángulo ($e_1 \times H$) más un triángulo
(ensanchamiento $(e_2 - e_1) \times H / 2$):

\begin{align*}
  A_\text{rect} &= e_1 \cdot H = 0{,}60 \times 7{,}30 = 4{,}38~\text{m}^2 \\
  A_\text{tri}  &= \frac{e_2 - e_1}{2} \cdot H
                = \frac{0{,}75 - 0{,}60}{2} \times 7{,}30 = 0{,}548~\text{m}^2 \\[4pt]
  W_f &= \gamma_\text{HA} \cdot (A_\text{rect} + A_\text{tri})
       = 2{,}50 \times (4{,}38 + 0{,}548)
       = \resultado{12{,}32~\text{Tn/ml}}
\end{align*}

Para obtener el brazo $\bar{x}_f$ se pondera cada subcomponente:

\begin{align*}
  W_{f,\text{rect}} &= 2{,}50 \times 4{,}38 = 10{,}95~\text{Tn/ml};
    \quad x_1 = B_\text{punta} + \frac{e_1}{2}
             = 0{,}80 + \frac{0{,}60}{2} = 1{,}10~\text{m} \\
  W_{f,\text{tri}}  &= 2{,}50 \times 0{,}548 = 1{,}37~\text{Tn/ml};
    \quad x_2 = B_\text{punta} + \frac{2\,(e_2 - e_1)}{3}
             = 0{,}80 + \frac{2 \times 0{,}15}{3} = 0{,}90~\text{m} \\[6pt]
  \bar{x}_f &= \frac{10{,}95 \times 1{,}10 + 1{,}37 \times 0{,}90}{12{,}32}
             = \frac{12{,}045 + 1{,}233}{12{,}32} = 1{,}07~\text{m} \\[4pt]
  M_f &= 1{,}07 \times 12{,}32 = \resultado{13{,}18~\text{Tn·m/ml}}
\end{align*}

\paragraph{C) Tierra no saturada ($h_{t,ns} = 2{,}50$~m, $\gamma_{t,ns} = 1{,}82$~Tn/m³):}

\begin{align*}
  N_{t,ns} &= \gamma_{t,ns} \cdot B_\text{talón} \cdot h_{t,ns}
            = 1{,}82 \times 2{,}45 \times 2{,}50
            = \resultado{11{,}15~\text{Tn/ml}} \\
  M_{t,ns} &= x_\text{talón} \cdot N_{t,ns}
            = 2{,}78 \times 11{,}15 = \resultado{31{,}00~\text{Tn·m/ml}}
\end{align*}

\paragraph{D) Tierra saturada ($h_{t,s} = 7{,}30 - 2{,}50 = 4{,}80$~m, $\gamma_{t,s} = 1{,}98$~Tn/m³):}

\begin{align*}
  N_{t,s} &= \gamma_{t,s} \cdot B_\text{talón} \cdot h_{t,s}
           = 1{,}98 \times 2{,}45 \times 4{,}80
           = \resultado{23{,}30~\text{Tn/ml}} \\
  M_{t,s} &= 2{,}78 \times 23{,}30 = \resultado{64{,}77~\text{Tn·m/ml}}
\end{align*}

\paragraph{E) Sobrecarga de tráfico ($q = 1$~Tn/m²):}

\begin{align*}
  W_q &= q \cdot B_\text{talón} = 1{,}00 \times 2{,}45
       = \resultado{2{,}45~\text{Tn/ml}} \\
  M_q &= 2{,}78 \times 2{,}45 = \resultado{6{,}81~\text{Tn·m/ml}}
\end{align*}

\begin{table}[H]
  \centering
  \caption{Fuerzas verticales y momentos estabilizadores — $B = 4{,}00$~m (sin subpresión)}
  \label{tab:fuerzas_verticales_4m}
  \begin{tabular}{@{}lccc@{}}
    \toprule
    \textbf{Elemento} & $N$ \textbf{(Tn/ml)} & $x$ \textbf{(m)} & $M_s$ \textbf{(Tn·m/ml)} \\
    \midrule
    Zapata ($4{,}00 \times 0{,}80$~m)          & 8{,}00  & 2{,}00 & 16{,}00 \\
    Fuste (trapezoidal)                          & 12{,}32 & 1{,}07 & 13{,}18 \\
    Tierra no saturada ($h = 2{,}50$~m)         & 11{,}15 & 2{,}78 & 31{,}00 \\
    Tierra saturada ($h = 4{,}80$~m)            & 23{,}30 & 2{,}78 & 64{,}77 \\
    Sobrecarga tráfico ($q = 1$~Tn/m²)          & 2{,}45  & 2{,}78 &  6{,}81 \\
    \midrule
    \textbf{TOTAL} & $\mathbf{\Sigma N = 57{,}22}$ & — & $\mathbf{\Sigma M_s = 131{,}76}$ \\
    \bottomrule
  \end{tabular}
\end{table}

\begin{notabox}[Subpresión en el primer tanteo]
  En este primer tanteo ($B = 4{,}00$~m) no se incluye la subpresión
  hidrostática bajo la zapata. La comprobación al vuelco ya resulta
  desfavorable sin ella ($FS = 1{,}76 < 1{,}80$), por lo que no es
  necesario añadirla para tomar la decisión de ampliar la base. La
  subpresión se incorpora explícitamente en el recálculo con
  $B = 5{,}60$~m (apartado~\ref{subsec:fuerzas_verticales_56m}), donde
  conduce a un resultado más conservador.
\end{notabox}

% -----------------------------------------------------------
\subsection{Empujes activos de Rankine}
\label{subsec:empujes_rankine_opc}
% -----------------------------------------------------------

Se aplica la \textbf{teoría de Rankine} para relleno horizontal y trasdós
vertical. El coeficiente de empuje activo vale:

\begin{equation}
  K_a = \tan^2\!\left(45^\circ - \frac{\phi'}{2}\right)
  \label{eq:ka_opc}
\end{equation}

\paragraph{Coeficientes por estrato:}
\begin{align*}
  \text{Arena limosa } (\phi' = 29^\circ): &\quad
    K_{a,1} = \tan^2(30{,}5^\circ) = \resultado{0{,}347} \\
  \text{Arcillas versicolores } (\phi' = 25{,}2^\circ): &\quad
    K_{a,2} = \tan^2(32{,}4^\circ) = \resultado{0{,}402}
\end{align*}

La presión horizontal efectiva en el estrato arcilloso incluye la cohesión:
\begin{equation}
  \sigma'_{h,a} = K_{a,2}\,\sigma'_v - 2\,c'\sqrt{K_{a,2}}
  \label{eq:sigma_ha_coh}
\end{equation}

En la separación de presiones se trata siempre por separado la presión efectiva
del suelo y la presión intersticial del agua.

\subsubsection{Estrato 1 — Arena limosa no saturada ($h_1 = 2{,}50$~m)}

\begin{align*}
  e_1 &= K_{a,1} \cdot \gamma_\text{nat} \cdot h_1
        = 0{,}347 \times 1{,}82 \times 2{,}50 = 1{,}579~\text{Tn/m}^2 \\[4pt]
  E_1 &= \tfrac{1}{2}\,e_1\,h_1
        = \tfrac{1}{2} \times 1{,}579 \times 2{,}50 = \resultado{1{,}97~\text{Tn/ml}} \\
  x_1 &= h_z + (H - h_1) + \frac{h_1}{3}
        = 0{,}80 + 4{,}80 + \frac{2{,}50}{3} = 6{,}43~\text{m} \\
  M_1 &= E_1 \cdot x_1 = 1{,}97 \times 6{,}43 = \resultado{12{,}67~\text{Tn·m/ml}}
\end{align*}

\begin{notabox}[Tensiones verticales en cotas clave]
  Para el cálculo de los empujes se emplean las tensiones verticales efectivas
  obtenidas del sondeo:
  \begin{itemize}
    \item Cota $-2{,}50$~m (techo N.F.):
      $\sigma_v = 1{,}82 \times 2{,}50 = 4{,}55$~Tn/m²;~$u=0$;
      $\sigma'_v = 4{,}55$~Tn/m²
    \item Cota $-3{,}50$~m (base arena saturada):
      $\sigma_v = 4{,}55 + 1{,}98 \times 1{,}00 = 6{,}53$~Tn/m²;
      $u = 1{,}00$~Tn/m²; $\sigma'_v = 5{,}53$~Tn/m²
  \end{itemize}
\end{notabox}

\subsubsection{Estrato 1b — Arena limosa saturada ($h_2 = 1{,}00$~m)}

En este estrato la presión efectiva horizontal varía de $e_1 = 1{,}579$~Tn/m² (en la
parte superior, al entrar el agua) a:
\begin{align*}
  e_2 &= K_{a,1}\!\left(\gamma_\text{nat}\,h_1
        + (\gamma_\text{sat} - \gamma_w)\,h_2\right)
       = 0{,}347\,(1{,}82 \times 2{,}50 + 0{,}98 \times 1{,}00)
       = 0{,}347 \times 5{,}53 = 1{,}92~\text{Tn/m}^2
\end{align*}

El empuje se descompone en un bloque rectangular y un triángulo:
\begin{align*}
  E_2 &= e_1 \cdot h_2 + \tfrac{1}{2}(e_2 - e_1)\,h_2
        = 1{,}579 \times 1{,}00 + \tfrac{1}{2}(1{,}92 - 1{,}579)\times 1{,}00
        = \resultado{1{,}75~\text{Tn/ml}}
\end{align*}

El brazo de aplicación se obtiene ponderando los momentos de cada sub-bloque:
\begin{align*}
  E_\text{cuad} &= e_1 \cdot h_2 = 1{,}579~\text{Tn/ml};
    \quad \text{brazo} = h_z + h_3 + h_2/2 = 0{,}80 + 3{,}80 + 0{,}50 = 5{,}10~\text{m}\\
  E_\text{tri}  &= \tfrac{1}{2}(e_2-e_1)\,h_2 = 0{,}171~\text{Tn/ml};
    \quad \text{brazo} = h_z + h_3 + h_2/3 = 0{,}80 + 3{,}80 + 0{,}33 = 4{,}93~\text{m}\\[4pt]
  x_2 &= \frac{1{,}579\times5{,}10 + 0{,}171\times4{,}93}{1{,}75}
        = \frac{8{,}05 + 0{,}84}{1{,}75} = 5{,}08~\text{m}\\
  M_2 &= 1{,}75 \times 5{,}08 = \resultado{8{,}89~\text{Tn·m/ml}}
\end{align*}

\subsubsection{Empuje del agua ($h_w = 4{,}80 + 0{,}80 = 5{,}60$~m)}

\begin{align*}
  E_w &= \tfrac{1}{2}\,\gamma_w\,h_w^2
        = \tfrac{1}{2} \times 1{,}00 \times (4{,}80+0{,}80)^2
        = \resultado{15{,}68~\text{Tn/ml}} \\
  M_w &= E_w \cdot \frac{h_w}{3}
        = 15{,}68 \times \frac{4{,}80+0{,}80}{3}
        = \resultado{29{,}27~\text{Tn·m/ml}}
\end{align*}

\subsubsection{Estrato 2 — Arcillas versicolores ($h_3 = 3{,}80$~m)}

Se comprueba la presión efectiva horizontal en la base del muro (cota $-7{,}30$~m):
\begin{align*}
  \sigma_{v_3} &= 6{,}53 + \gamma_3\,h_3 = 6{,}53 + 1{,}78 \times (8{,}00-3{,}50)
                = 14{,}54~\text{Tn/m}^2 \\
  u_3 &= \gamma_w \cdot (8{,}00 - 2{,}50) = 5{,}50~\text{Tn/m}^2 \\
  \sigma'_{v_3} &= 14{,}54 - 5{,}50 = 9{,}04~\text{Tn/m}^2 \\[6pt]
  \sigma_{H_3} &= K_{a,2}\,\sigma'_{v_3} - 2\,c'\sqrt{K_{a,2}}
    = 0{,}402 \times 9{,}04 - 2 \times 4{,}70 \times \sqrt{0{,}402}
    = 3{,}63 - 5{,}96 = -2{,}15~\text{Tn/m}^2 \quad 
\end{align*}

Como $\sigma_{H_3} < 0$, la cohesión suprime el empuje activo del suelo en toda la
zona de arcillas: el terreno no tracciona sobre el muro. Sin embargo, la presión
horizontal en la parte superior de la arcilla (heredada del estrato anterior) es
positiva ($e_2 = 1{,}92$~Tn/m²) y en la base nula, obteniéndose un diagrama
triangular:

\begin{align*}
  E_3 &= \tfrac{1}{2}\,e_2\,h_3
        = \tfrac{1}{2} \times 1{,}92 \times 3{,}80 = \resultado{3{,}64~\text{Tn/ml}} \\
  M_3 &= E_3 \cdot \left(h_z + \frac{2\,h_3}{3}\right)
        = 3{,}64 \times \left(0{,}80 + \frac{2\times3{,}80}{3}\right)
        = 3{,}64 \times 3{,}33 = \resultado{12{,}13~\text{Tn·m/ml}}
\end{align*}

\subsubsection{Sobrecarga de tráfico ($q = 1$~Tn/m²)}

La sobrecarga genera un empuje horizontal equivalente en cada estrato:

\paragraph{En la zona de arena ($h_1 + h_2 = 3{,}50$~m):}
\begin{align*}
  E_{q,1} &= q \cdot K_{a,1} \cdot (h_1+h_2)
            = 1{,}00 \times 0{,}347 \times 3{,}50 = \resultado{1{,}21~\text{Tn/ml}} \\
  x_{q,1} &= h_z + h_3 + \frac{h_1+h_2}{2}
            = 0{,}80 + 3{,}80 + \frac{3{,}50}{2} = 6{,}35~\text{m} \\
  M_{q,1} &= 1{,}21 \times 6{,}35 = \resultado{7{,}68~\text{Tn·m/ml}}
\end{align*}

\paragraph{En la zona de arcilla ($h_3 = 3{,}80$~m):}
\begin{align*}
  E_{q,2} &= q \cdot K_{a,2} \cdot h_3
            = 1{,}00 \times 0{,}402 \times 3{,}80 = \resultado{1{,}53~\text{Tn/ml}} \\
  x_{q,2} &= h_z + \frac{h_3}{2} = 0{,}80 + \frac{3{,}80}{2} = 2{,}70~\text{m} \\
  M_{q,2} &= 1{,}53 \times 2{,}70 = \resultado{4{,}13~\text{Tn·m/ml}}
\end{align*}

\subsubsection{Resumen del empuje activo total}

\begin{table}[H]
  \centering
  \caption{Integración del diagrama de empuje activo (por ml de muro)}
  \label{tab:empuje_activo_opc}
  \begin{tabular}{@{}lcccc@{}}
    \toprule
    \textbf{Componente} & \textbf{Forma}
    & $E$ \textbf{(Tn/ml)} & $x$ \textbf{(m)} & $M$ \textbf{(Tn·m/ml)} \\
    \midrule
    $E_1$ — Arena no sat.\ (triáng.) & Triáng.\ & 1{,}97 & 6{,}43 & 12{,}67 \\
    $E_2$ — Arena sat.\ (rect.+triáng.) & Mixta  & 1{,}75 & 5{,}08 &  8{,}89 \\
    $E_w$ — Agua ($h_w=5{,}60$~m) & Triáng.\  & 15{,}68 & 1{,}87 & 29{,}27 \\
    $E_3$ — Arcilla (triáng.\ decreciente) & Triáng.\ & 3{,}64 & 3{,}33 & 12{,}13 \\
    $E_{q,1}$ — Sobrecarga en arena & Rect.\ & 1{,}21 & 6{,}35 & 7{,}68 \\
    $E_{q,2}$ — Sobrecarga en arcilla & Rect.\  & 1{,}53 & 2{,}70 & 4{,}13 \\
    \midrule
    \textbf{TOTAL} & & $\mathbf{E_T = 25{,}78}$ & — & $\mathbf{M_T = 74{,}77}$ \\
    \bottomrule
  \end{tabular}
\end{table}

\begin{resultadobox}[Empuje activo total]
  \[
    E_T = 25{,}78~\text{Tn/ml}
    \qquad
    M_T = 74{,}77~\text{Tn·m/ml}
  \]
  Altura de aplicación de la resultante sobre la base de la zapata:
  \[
    d_T = \frac{M_T}{E_T} = \frac{74{,}77}{25{,}78} = 2{,}90~\text{m}
  \]
\end{resultadobox}

% -----------------------------------------------------------
\subsection{Comprobación al vuelco — Primer tanteo ($B = 4{,}00$~m)}
\label{subsec:vuelco_4m}
% -----------------------------------------------------------

El factor de seguridad al vuelco respecto a la punta de la zapata es:

\begin{equation}
  FS_\text{vuelco} = \frac{\sum M_s}{\sum M_v}
    = \frac{131{,}76}{74{,}77} = \resultado{1{,}76 < 1{,}80 \quad \text{\texttimes\ NO CUMPLE}}
  \label{eq:fs_vuelco_4m}
\end{equation}

\begin{supuestobox}[Decisión: ampliar la base de zapata a B = 5 m]
  Con $B = 4{,}00$~m el factor de seguridad al vuelco ($FS = 1{,}76$) no alcanza
  el mínimo exigido ($FS \geq 1{,}80$). Dado que el momento volcador es fijo
  (depende exclusivamente del empuje), la única solución es aumentar el momento
  estabilizador ampliando el \textbf{talón} de la zapata.

  Se adopta \textbf{$B = 5{,}60$~m}, lo que amplía el talón a:
  \[
    B_{\text{talón}} = 5{,}60 - 0{,}75 - 0{,}80 = 4{,}05~\text{m}
  \]
\end{supuestobox}

% -----------------------------------------------------------
\subsection{Fuerzas verticales y subpresión — $B = 5{,}60$~m}
\label{subsec:fuerzas_verticales_56m}
% -----------------------------------------------------------

De la misma forma que en apartado 2.1.2, se calcula el peso por metro lineal de cada componente y su distancia a la
punta ($x = 0$ en la punta, positivo hacia el trasdós).
Con $B_\text{talón} = 4{,}05$~m, el brazo de las cargas del talón es:
\[
  x_\text{talón} = B_\text{punta} + e_2 + \frac{B_\text{talón}}{2}
    = 0{,}80 + 0{,}75 + \frac{4{,}05}{2} = \mathbf{3{,}58}~\text{m}
\]

\paragraph{Subpresión bajo la zapata.}

El nivel freático a $-2{,}50$~m genera una presión hidrostática ascendente en
la base de la zapata (cota $-8{,}10$~m). La distribución es \textbf{triangular}:
máxima en el trasdós (donde $h_w = 8{,}10 - 2{,}50 = 5{,}60$~m) y nula en la
puntera, gracias al drenaje del intradós.

\begin{align*}
  N_{s,w} &= \frac{1}{2}\,\gamma_w\,h_w\,B
            = \frac{1}{2} \times 1{,}00 \times 5{,}60 \times 5{,}60
            = \resultado{15{,}68~\text{Tn/ml}} \quad (\uparrow) \\[4pt]
  x_{s,w} &= \frac{2}{3}\,B = \frac{2}{3} \times 5{,}60 = 3{,}73~\text{m}
            \quad \text{(resultante del triángulo, desde la punta)} \\[4pt]
  M_{s,w} &= N_{s,w} \times x_{s,w} = 15{,}68 \times 3{,}73
            = \resultado{58{,}54~\text{Tn·m/ml}} \quad (\text{reduce}~\Sigma M_s)
\end{align*}

\begin{table}[H]
  \centering
  \caption{Fuerzas verticales, subpresión y momentos estabilizadores — $B = 5{,}60$~m}
  \label{tab:fuerzas_verticales_5m}
  \begin{tabular}{@{}lcccc@{}}
    \toprule
    \textbf{Elemento} & $N$ \textbf{(Tn/ml)} & $x$ \textbf{(m)} & $M_s$ \textbf{(Tn·m/ml)} \\
    \midrule
    Zapata ($5{,}60 \times 0{,}80$)       & $+$11{,}20 & 2{,}80 & $+$31{,}36 \\
    Fuste (trapezoidal)                    & $+$12{,}32 & 1{,}07 & $+$13{,}18 \\
    Tierra no saturada ($h=2{,}50$~m)     & $+$18{,}45 & 3{,}58 & $+$65{,}96 \\
    Tierra saturada ($h=4{,}80$~m)        & $+$38{,}49 & 3{,}58 & $+$137{,}60 \\
    Sobrecarga tráfico ($q=1$~Tn/m²)      & $+$4{,}05  & 3{,}58 & $+$14{,}48 \\
    \midrule
    Subtotal (sin subpresión)              & 84{,}51    & —      & 262{,}58 \\
    \midrule
    ($-$)~Subpresión $N_{s,w}$ ($\uparrow$) & $-$15{,}68 & 3{,}73 & $-$58{,}54 \\
    \midrule
    \textbf{TOTAL} & $\mathbf{N_T = 68{,}80}$ & — & $\mathbf{M_s = 203{,}96}$ \\
    \bottomrule
  \end{tabular}
\end{table}

\begin{notabox}[Verificación de fuerzas]
  \begin{itemize}
    \item Zapata: $N = 2{,}50 \times 5{,}60 \times 0{,}80 = 11{,}20$~Tn/ml;\quad
          $M = 11{,}20 \times 2{,}80 = 31{,}36$~Tn·m/ml
    \item Tierra no sat.: $N = 1{,}82 \times 2{,}50 \times 4{,}05 = 18{,}45$~Tn/ml;\quad
          $M = 18{,}45 \times 3{,}58 = 65{,}96$~Tn·m/ml
    \item Tierra sat.: $N = 1{,}98 \times 4{,}80 \times 4{,}05 = 38{,}49$~Tn/ml;\quad
          $M = 38{,}49 \times 3{,}58 = 137{,}60$~Tn·m/ml
    \item Sobrecarga: $N = 1{,}00 \times 4{,}05 = 4{,}05$~Tn/ml;\quad
          $M = 4{,}05 \times 3{,}58 = 14{,}48$~Tn·m/ml
    \item Subpresión triangular: $N_{s,w} = \tfrac{1}{2} \times 5{,}60 \times 5{,}60 = 15{,}68$~Tn/ml;
          brazo $= \tfrac{2}{3} \times 5{,}60 = 3{,}73$~m;\quad $M = 58{,}54$~Tn·m/ml
  \end{itemize}
\end{notabox}

% -----------------------------------------------------------
\subsection{Comprobación al vuelco ($B = 5{,}60$~m)}
\label{subsec:vuelco_5m}
% -----------------------------------------------------------

\begin{equation}
  FS_\text{vuelco} = \frac{\sum M_s}{\sum M_v}
    = \frac{203{,}96}{74{,}77} = \resultado{2{,}73 \geq 1{,}80 \quad \checkmark~\text{CUMPLE}}
  \label{eq:fs_vuelco_5m}
\end{equation}

% -----------------------------------------------------------
\subsection{Comprobación al deslizamiento}
\label{subsec:deslizamiento}
% -----------------------------------------------------------

La resistencia al deslizamiento en la base de la zapata se evalúa considerando
el contacto con el \textbf{estrato~3} (arcillas margosas,
$\phi' = 26^\circ$, $c' = 7{,}80$~Tn/m²). Se incluyen la fuerza de rozamiento,
la cohesión, el empuje pasivo del suelo y el empuje pasivo del agua en el
intradós.

\paragraph{Empuje pasivo del suelo ($D_f = 1{,}00$~m):}

\begin{equation}
  K_{p,2} = \tan^2\!\left(45^\circ + \frac{25{,}2^\circ}{2}\right) = 2{,}49
\end{equation}

\begin{align*}
  E_p &= \tfrac{1}{2}\,K_{p,2}\,(\gamma_3-\gamma_w)\,D_f^2
        = \tfrac{1}{2} \times 2{,}49 \times (1{,}78-1{,}00) \times 1{,}00^2
        = \resultado{0{,}97~\text{Tn/ml}} \\
  M_p &= D_f \cdot \frac{1}{3} \cdot E_p = 1{,}00 \times \frac{1}{3} \times 0{,}97
        = \resultado{0{,}32~\text{Tn·m/ml}}
\end{align*}

\paragraph{Empuje pasivo del agua:}

\begin{align*}
  E_{p,w} &= \tfrac{1}{2}\,\gamma_w\,h_w^2
            = \tfrac{1}{2} \times 1{,}00 \times 1{,}00^2 = \resultado{0{,}50~\text{Tn/ml}} \\
  M_{p,w} &= D_f \cdot \frac{1}{3} \cdot E_{p,w}
            = 1{,}00 \times \frac{1}{3} \times 0{,}50 = \resultado{0{,}17~\text{Tn·m/ml}}
\end{align*}

\paragraph{Fuerza de rozamiento en la base:}

El ángulo de rozamiento zapata-terreno se adopta como $\delta = \tfrac{2}{3}\,\phi'$:
\begin{align*}
  \mu &= \tan\!\left(\frac{2}{3}\,\phi'\right)
        = \tan\!\left(\frac{2}{3} \times 26^\circ\right)
        = \tan(17{,}33^\circ) = 0{,}31 \\[4pt]
  F   &= \mu \cdot N_T + c' \cdot A
        = 0{,}31 \times 68{,}83 + 7{,}80 \times 5{,}60
        = 21{,}34 + 43{,}68 = 65{,}02~\text{Tn}
\end{align*}

\paragraph{Fuerza estabilizadora total:}
\[
  F_e = F + E_p + E_{p,w} = 65{,}02 + 0{,}97 + 0{,}50 = 66{,}48~\text{Tn}
\]

\begin{equation}
  C_{sd} = \frac{F_e}{E_T} = \frac{66{,}48}{25{,}78} = \resultado{2{,}58 \geq 1{,}50 \quad \checkmark~\text{CUMPLE}}
  \label{eq:fs_desliz}
\end{equation}

% -----------------------------------------------------------
\subsection{Comprobación al hundimiento}
\label{subsec:hundimiento}
% -----------------------------------------------------------

Se adopta una \textbf{tensión admisible} $\sigma_\text{adm} = 2$~kg/cm²
$= 20$~Tn/m².

\paragraph{Posición de la resultante vertical:}

El momento total neto respecto a la punta es:
\begin{align*}
  M_T &= \Sigma M_s + M_{p,\text{pasivo}} - \Sigma M_v
        = 204{,}04 + (0{,}32 + 0{,}17) - 74{,}77 = 129{,}76~\text{Tn·m/ml}
\end{align*}

La resultante vertical actúa a una distancia $x_R$ desde la punta:
\[
  x_R = \frac{M_T}{N_T} = \frac{129{,}76}{68{,}83} = 1{,}88~\text{m}
\]

Comprobación del núcleo central (tercio medio):
\[
  \frac{B}{3} = \frac{5{,}60}{3} = 1{,}87~\text{m} \leq x_R = 1{,}88~\text{m}
  \leq \frac{2B}{3} = 3{,}73~\text{m} \quad \checkmark \quad
  \text{La resultante está dentro del núcleo central}
\]

\paragraph{Distribución de presiones en la base (trapezoidal):}

\begin{align*}
  \sigma &= \frac{N_T}{B} \pm \frac{6\,M_T}{B^2}
           = \frac{68{,}83}{5{,}60} \pm \frac{6 \times 129{,}76}{5{,}60^2}
           = 12{,}29 \pm 24{,}82~\text{Tn/m}^2 \\[4pt]
  \sigma_1 &= 12{,}29 + 24{,}82 = 37{,}11~\text{Tn/m}^2 = \resultado{3{,}71~\text{kg/cm}^2} \\
  \sigma_2 &= 12{,}29 - 24{,}82 = -12{,}53~\text{Tn/m}^2 = -1{,}25~\text{kg/cm}^2
\end{align*}

\paragraph{Comprobaciones:}

\begin{enumerate}
  \item $\sigma_1 \geq 1{,}25\,\sigma_\text{adm} = 2{,}50$~kg/cm²
    $\longrightarrow$ $3{,}71 \geq 2{,}50$ \checkmark~\textbf{CUMPLE}
  \item $\dfrac{\sigma_1 + \sigma_2}{2} \leq \sigma_\text{adm}$
    $\longrightarrow$
    $\dfrac{3{,}71 - 1{,}25}{2} = 1{,}23 \leq 2{,}00$ \checkmark~\textbf{CUMPLE}
\end{enumerate}

\begin{table}[H]
  \centering
  \caption{Verificación de estabilidad — muro en ménsula ($B = 5{,}60$~m, con subpresión)}
  \label{tab:estabilidad_muro_opc}
  \begin{tabular}{@{}lcccc@{}}
    \toprule
    \textbf{Comprobación} & \textbf{FS / Valor}
    & \textbf{Límite} & \textbf{Estado} \\
    \midrule
    Vuelco ($B=4{,}00$~m, ref.) & 1{,}76 & $\geq 1{,}80$ & \textbf{NO CUMPLE} \\
    \midrule
    Vuelco ($B=5{,}60$~m) & 2{,}73 & $\geq 1{,}80$ & \checkmark \textbf{CUMPLE} \\
    Deslizamiento & 1{,}88 & $\geq 1{,}50$ & \checkmark \textbf{CUMPLE} \\
    Núcleo central ($x_R$) & 1{,}88~m & $\geq B/3 = 1{,}87$~m & \checkmark \textbf{CUMPLE} \\
    Hundimiento ($\sigma_\text{med}$) & 1{,}23~kg/cm² & $\leq 2{,}00$ & \checkmark \textbf{CUMPLE} \\
    \bottomrule
  \end{tabular}
\end{table}

% -----------------------------------------------------------
\subsubsection{Proceso constructivo del muro en ménsula}
% -----------------------------------------------------------

\begin{enumerate}
  \item Excavación del intradós hasta la cota de la base de zapata ($D_f=1{,}00$~m
        bajo rasante intradós), con entibación ligera si procede.
  \item Ejecución del \textbf{hormigón de limpieza} HM-15 ($e=10$~cm) sobre el
        fondo de excavación.
  \item Colocación de la \textbf{armadura} de la zapata y del arranque del fuste.
  \item Hormigonado de la zapata con HA-30/B/20/IIb+QA con cemento SR.
  \item Encofrado y hormigonado del \textbf{fuste} por tongadas de \mbox{$\sim$3~m},
        con vibrado y curado.
  \item Instalación del sistema de \textbf{drenaje} del trasdós: lámina
        geotextil filtrante + capa de grava ($\varnothing$~20--40~mm) + tubo
        dren $\varnothing\,100$~mm en la base del trasdós.
  \item \textbf{Relleno del trasdós} por tongadas de 30~cm, compactando con
        rodillo vibrante hasta obtener el 95\% del Proctor Modificado.
  \item Juntas de dilatación cada 15--20~m (junta de poliestireno expandido
        $e=20$~mm con sellante de poliuretano).
\end{enumerate}

% -----------------------------------------------------------
\subsection{Asientos del muro en ménsula ($B = 5{,}60$~m)}
\label{subsec:asientos_muro_56m}
% -----------------------------------------------------------

La zapata queda a cota $-8{,}10$~m, apoyada sobre el \textbf{estrato~3}
(arcillas margosas, $-8$ a $-25$~m, $H = 17$~m). La tensión inicial
efectiva en la cota de la base de la zapata (subpresión incluida):
\[
  \sigma'_{v,\text{zapata}} = 9{,}04 + (1{,}92-1)\times 0{,}10 = 9{,}13~\text{Tn/m}^2
\]

Carga neta transmitida al terreno:
\[
  \sigma_\text{zapata} = \frac{N_T}{B} = \frac{68{,}80}{5{,}60} = 12{,}29~\text{Tn/m}^2
  \qquad
  \Delta\sigma = 12{,}29 - 9{,}13 = \resultado{3{,}16~\text{Tn/m}^2}
\]

\subsubsection{Método 1 — Ensayo edométrico (por subcapas)}

Se adoptan los parámetros edométricos del estrato~3: $C_c = 0{,}15$,
$e_0 = 0{,}60$. El estrato se divide en dos subcapas de $H = 8{,}5$~m
para mayor precisión.

\paragraph{Subcapa A ($-8$ a $-16{,}5$~m, $H_{A} = 8{,}5$~m):}
\begin{align*}
  \sigma'_{i,3a} &= 9{,}04 + (1{,}92-1)\times\frac{8{,}5}{2} = 12{,}95~\text{Tn/m}^2
    \quad\text{(punto representativo = } H/2\text{)} \\
  \sigma'_{f,3a} &= 12{,}95 + 3{,}16 = 16{,}11~\text{Tn/m}^2 \\
  \Delta e_{3a}  &= 0{,}15\cdot\log\!\left(\frac{16{,}11}{12{,}95}\right) = 0{,}01422 \\
  \varepsilon_{3a} &= \frac{\Delta e_{3a}}{1+e_0} = \frac{0{,}01422}{1{,}60} = 0{,}00889 \\
  S_{3a} &= \varepsilon_{3a}\times H_A = 0{,}00889\times 8{,}5 = \resultado{7{,}6~\text{cm}}
\end{align*}

\paragraph{Subcapa B ($-16{,}5$ a $-25$~m, $H_{B} = 8{,}5$~m):}
\begin{align*}
  \sigma'_{i,3b} &= 9{,}04 + (1{,}92-1)\times\!\left(8{,}5+\frac{8{,}5}{2}\right)
                  = 20{,}77~\text{Tn/m}^2 \\
  \sigma'_{f,3b} &= 20{,}77 + 3{,}16 = 23{,}93~\text{Tn/m}^2 \\
  \Delta e_{3b}  &= 0{,}15\cdot\log\!\left(\frac{23{,}93}{20{,}77}\right) = 0{,}00922 \\
  \varepsilon_{3b} &= \frac{0{,}00922}{1{,}60} = 0{,}00576 \\
  S_{3b} &= 0{,}00576\times 8{,}5 = \resultado{4{,}9~\text{cm}}
\end{align*}

\begin{equation}
  \resultado{S_{\text{total, edométrico}} = 7{,}6 + 4{,}9 = 12{,}5~\text{cm}}
  \label{eq:asiento_edom}
\end{equation}

\subsubsection{Método 2 — Steinbrenner (asiento elástico + consolidación)}

Datos del estrato~3: $E_3 = 25{,}278$~MPa $= 2\,527{,}8$~Tn/m²,
$\nu_3 = 0{,}25$. Se trabaja por metro lineal ($n = L/B \to \infty$;
se toma $n=10$, $\phi_1 = 2{,}322$).

\paragraph{Asiento inmediato:}
\begin{align*}
  A &= 1 - \nu_3^2 = 1 - 0{,}25^2 = 0{,}9375 \\
  S_i &= \frac{\Delta p_{\text{net}}\cdot B}{2\,E_3}\,(A\cdot\phi_1)
       = \frac{12{,}29\times 5{,}60}{2\times 2\,527{,}8}
         \times(0{,}9375\times 2{,}322)
       = \resultado{2{,}97~\text{cm}}
\end{align*}

\paragraph{Asiento de consolidación ($H_{\mathrm{ef}} = 2B = 11{,}20$~m):}
\begin{align*}
  E_{ed} &= E_3\,\frac{1-\nu_3}{(1+\nu_3)(1-2\nu_3)}
           = 2\,527{,}8\times\frac{0{,}75}{1{,}25\times 0{,}50}
           = 3\,033{,}4~\text{Tn/m}^2 \\
  S_c    &= \frac{\Delta p_{\text{net}}}{E_{ed}}\cdot H_{\mathrm{ef}}
           = \frac{12{,}29}{3\,033{,}4}\times 11{,}20
           = \resultado{4{,}54~\text{cm}}
\end{align*}

\begin{equation}
  S_t = S_i + S_c = 2{,}97 + 4{,}54 = \resultado{7{,}51~\text{cm}}
  \label{eq:asiento_steinbrenner}
\end{equation}

\begin{resultadobox}[Asiento total por Steinbrenner]
  \[
    S_t = 7{,}51~\text{cm} > 5{,}0~\text{cm}
    \quad\Longrightarrow\quad
    \textbf{DEBEMOS PILOTAR la estructura}
  \]
\end{resultadobox}

\begin{notabox}[Comparativa de métodos]
  El método edométrico arroja $S = 12{,}5$~cm y Steinbrenner
  $S = 7{,}5$~cm; ambos superan el límite admisible de $5$~cm para
  muros de HA. Además, aumentar el ancho de la zapata agravaría el
  asiento al incrementar la carga transmitida al estrato~3.
\end{notabox}

% -----------------------------------------------------------
\subsection{Estimación de pilotes}
\label{subsec:pilotes_muro}
% -----------------------------------------------------------

Se adopta un \textbf{pilote aislado por metro lineal} que recibe la
totalidad de la carga: $P = N_T = 68{,}80$~Tn/ml (hormigonado in situ,
empotrado en el estrato~3).

\paragraph{Datos del pilote:}
$D = 0{,}45$~m, $A_p = 0{,}159$~m²;
$E_p = 30\,000$~MPa $= 3\,060\,000$~Tn/m²;
$N_\mathrm{SPT} = 33$ (estrato~3); $f_N = 0{,}20$; $f_f = 2$;
$\gamma_R = 1{,}5$.

\paragraph{Resistencia por punta:}
\begin{equation}
  R_{pk} = f_N\cdot N_\mathrm{SPT}\cdot 1\,000\cdot A_p
         = 0{,}20\times 33\times 1\,000\times 0{,}159
         = 1\,049~\text{kN} = \resultado{107~\text{Tn}}
\end{equation}

\paragraph{Resistencia por fuste ($L = 0{,}5$~m):}
\begin{align*}
  \tau_f &= f_f\cdot N_\mathrm{SPT} = 2\times 33 = 66~\text{kPa} \\
  A_f    &= \pi\cdot D\cdot L = \pi\times 0{,}45\times 0{,}5 = 0{,}707~\text{m}^2 \\
  R_{fk} &= \tau_f\cdot A_f = 66\times 0{,}707 = 46{,}6~\text{kN}
           = \resultado{4{,}76~\text{Tn}}
\end{align*}

\paragraph{Resistencia característica y de cálculo:}
\begin{align*}
  R_{ck} &= R_{pk} + R_{fk} = 107 + 4{,}76 = 111{,}76~\text{Tn} \\
  R_{cd} &= \frac{R_{ck}}{\gamma_R} = \frac{111{,}76}{1{,}5}
           = \resultado{74{,}5~\text{Tn} > 68{,}80~\text{Tn}} \quad\checkmark
\end{align*}

\paragraph{Asiento del pilote aislado:}

Con coeficiente de redistribución $\alpha$ intermedio ($x_p = 1$,
$\alpha_f = 0{,}5$):
\begin{align*}
  \alpha &= \frac{1}{R_{ck}}\!\left(0{,}5\,R_{fk} + R_{pk}\right)
           = \frac{0{,}5\times 4{,}76 + 107}{111{,}76} = 0{,}979 \\
  S_i    &= \left(\frac{D}{40\,R_{ck}} + \frac{\alpha\cdot L}{A_p\cdot E_p}\right)P
           = \left(\frac{0{,}45}{40\times 111{,}76}
             + \frac{0{,}979\times 0{,}5}{0{,}159\times 3\,060\,000}\right)
             \times 68{,}80
           = \resultado{0{,}44~\text{cm}}
\end{align*}

\paragraph{Asiento del grupo (Vesic):}
\begin{equation}
  S_g = \sqrt{\frac{B_g}{D}}\,S_i
      = \sqrt{\frac{4{,}0}{0{,}45}}\times 0{,}44
      = \resultado{0{,}66~\text{cm} < 2{,}5~\text{cm}} \quad\text{DESPRECIABLE}
\end{equation}

\begin{resultadobox}[Solución de cimentación adoptada — Muro]
  \textbf{1~pilote por metro lineal, $D = 0{,}45$~m, $L = 0{,}5$~m},
  hormigonado in situ, empotrado en el estrato~3 (arcillas margosas).
  El asiento del grupo $S_g = 0{,}66$~cm es despreciable gracias a
  la elevada rigidez del estrato~3. Tiene sentido: dicho estrato es
  tan rígido que cualquier pilote de pequeña longitud resuelve el
  asiento.
\end{resultadobox}

% -----------------------------------------------------------
\subsection{Opción B — Muro pantalla de hormigón armado}
\label{subsec:opcion_pantalla}
% -----------------------------------------------------------

\subsubsection{Descripción de la tipología}

El \textbf{muro pantalla} es una estructura de contención flexible ejecutada
\emph{in situ} antes de la excavación, sin necesidad de acceso por el lado
excavado. Su cualidad básica es la de contención flexible y la posibilidad
de actuar simultáneamente como \emph{cimentación profunda}. Los temas
principales de análisis son, según los criterios del curso: naturaleza del
terreno, rigidez a flexión, cálculo de empujes, profundidad de empotramiento
(\emph{clava}), influencia del agua, sustentación con anclajes y sistema de
ejecución.

Para $H = 7{,}30$~m se analiza la variante más adecuada:

\begin{description}
  \item[Pantalla libre (sin anclajes):] La pantalla actúa en voladizo,
  empotrada en el terreno. Apta para alturas $H \lesssim 4$--5~m; por
  encima de ese umbral los desplazamientos en coronación y los momentos
  flectores se vuelven prohibitivos.

  \item[Pantalla anclada (un nivel de anclajes):] Se instala un nivel de
  anclajes activos (\emph{micropilotes} o \emph{tirantes Gewi}) a
  $z_a \approx 1{,}0$~m bajo la coronación, reduciendo drásticamente el
  momento flector y la profundidad de empotramiento. Es la variante
  necesaria para $H > 5$~m.
\end{description}

\subsubsection{Proceso constructivo de la pantalla hormigonada \emph{in situ}}

De acuerdo con la teoría de clase (tema Pantallas, Villena~Manzanares):

\begin{enumerate}
  \item \textbf{Ejecución de muretes guía}: dos pequeños muros de hormigón
        en masa paralelos, separados el espesor de la pantalla más las
        tolerancias ($\sim$60--80~cm). Los muretes guían la cuchara bivalva
        y mantienen la verticalidad. Se deja un mínimo de 0{,}25~m por
        encima del nivel freático para evitar colapso del borde de la zanja.
  \item \textbf{Excavación de bataches}: se alternan bataches impares y pares
        (procedimiento de paneles alternados). Cada batache tiene longitud
        típica de 2{,}5--6~m. La excavación se realiza \textbf{con lodos
        bentoníticos} dado que el terreno es cohesivo--arenoso con N.F.\ a
        $-2{,}50$~m. El nivel de lodo debe mantenerse a \textbf{mínimo
        1{,}5~m sobre el nivel freático} (condición crítica del proceso).
  \item \textbf{Limpieza del fondo} y comprobación del contenido de arena
        en los lodos ($<\!6\%$ en peso).
  \item \textbf{Colocación de la jaula de armadura} mediante grúa.
  \item \textbf{Hormigonado por tubo Tremie} (hormigón HA-30 sumergido),
        desplazando el lodo bentonítico desde el fondo.
  \item \textbf{Extracción de elementos de junta} entre bataches.
  \item \textbf{Demolición de la cabeza} de los paneles hasta la cota de
        la viga de atado.
  \item \textbf{Ejecución de la viga de atado} (viga de coronación) que
        solidariza todos los paneles.
\end{enumerate}

\begin{notabox}[Condición de los lodos bentoníticos y sulfatos]
  La presencia de $3\,500$~mg~SO$_4^{2-}$/kg en el suelo condiciona la
  compatibilidad química de los lodos bentoníticos. Se deberá emplear
  bentonita con certificado de resistencia a sulfatos y verificar el
  pH de los lodos durante la ejecución ($7{,}5 < \text{pH} < 11{,}5$).
  El hormigón de la pantalla se especificará igualmente HA-30/SR.
\end{notabox}

\subsubsection{Espesor y profundidad de empotramiento}
\label{subsubsec:empotramiento}

\paragraph{Espesor adoptado:}
\[
  e_p \geq \frac{H}{10} = \frac{7{,}30}{10} = 0{,}73~\text{m}
  \quad\longrightarrow\quad \resultado{e_p = 0{,}75~\text{m}}
\]

\paragraph{Distancia de empotramiento real.}

La pantalla penetra a través del \textbf{estrato~2} (arcillas versicolores,
$\phi'=25{,}2^\circ$, $K_{p,2}=2{,}49$, espesor restante bajo la cota $-7{,}30$~m:
$0{,}70$~m) y continúa en el \textbf{estrato~3} (arcillas margosas,
$\phi'=26^\circ$, $K_{a,3}=0{,}390$, $K_{p,3}=2{,}561$). El empotramiento $t$ se
mide desde la base de la pantalla libre (cota $-7{,}30$~m).

Para la obtención de la distancia de empotramiento se igualan los momentos
activos y pasivos respecto a la base de la pantalla ($\Sigma M_a = \Sigma M_p$),
considerando las contribuciones de todos los estratos en activo (incluyendo
agua y sobrecarga) y el empuje pasivo en los estratos~2 y~3. Se pueden combrobar 
los cálculos en el anexo de cálulos manuales. Resuelta la
ecuación implícita en $t$ por equilibrio:

\begin{equation}
  \resultado{t = 14{,}68~\text{m}}
  \label{eq:empotramiento_pantalla}
\end{equation}

\begin{supuestobox}[Factor de seguridad estructural]
  Para tener en cuenta la simplificación del modelo (suelo homogéneo en
  cada estrato) se aplica un margen del 20\,\% sobre el empotramiento
  teórico:
  \[
    D = 1{,}20 \times 14{,}68 = \resultado{17{,}62~\text{m}}
  \]
  En un futuro se podría añadir \textbf{anclajes activos} para reducir
  la longitud de empotramiento necesaria.
\end{supuestobox}

La \textbf{longitud total de la pantalla} resulta:
\begin{equation}
  L_\text{pantalla} = H + D = 7{,}30 + 17{,}62 = \resultado{24{,}92~\text{m}}
  \label{eq:long_pantalla}
\end{equation}

\subsubsection{Comprobación frente a sifonamiento}
\label{subsubsec:sifonamiento}

Con el nivel freático en el trasdós a $-2{,}50$~m y el intradós drenado, la
diferencia de carga hidráulica es $\Delta H = h_w = 4{,}80$~m. La longitud
de percolación es la profundidad de empotramiento $D = 17{,}62$~m.

\paragraph{Gradiente crítico del estrato~3:}
\begin{equation}
  i_c = \frac{\gamma_{\mathrm{sum}}}{\gamma_w}
      = \frac{1{,}92 - 1{,}00}{1{,}00}
      = \resultado{0{,}92}
\end{equation}

\paragraph{Gradiente hidráulico real:}
\begin{equation}
  i = \frac{\Delta H}{D} = \frac{4{,}80}{17{,}62} = \resultado{0{,}272}
\end{equation}

\paragraph{Factor de seguridad frente a sifonamiento:}
\begin{equation}
  FS_{\mathrm{sifon}} = \frac{i_c}{i} = \frac{0{,}92}{0{,}272}
  = \resultado{3{,}38 > 3} \quad\checkmark \quad \textbf{CUMPLE}
  \label{eq:fs_sifonamiento}
\end{equation}

\paragraph{Distancia mínima de empotramiento por sifonamiento:}

Si en un futuro se añaden anclajes para reducir $D$, la condición
$FS \geq 3$ impone:
\[
  i_{\min} = \frac{i_c}{3} = \frac{0{,}92}{3} = 0{,}307
  \qquad\Rightarrow\qquad
  L_{\min} = \frac{\Delta H}{i_{\min}} = \frac{4{,}80}{0{,}307}
  = \resultado{15{,}65~\text{m}}
\]

Este resultado tiene una implicación práctica importante: si en el futuro se
añadiesen \textbf{anclajes activos} para reducir la profundidad de empotramiento
estructural (con el fin de abaratar el coste de la pantalla), la condición de
sifonamiento impondría un límite inferior de $L_\text{mín} = 15{,}65$~m que
\emph{no podría ser franqueado por los anclajes}. El motivo es que el sifonamiento
depende exclusivamente de la longitud del recorrido de percolación bajo la
pantalla, que coincide con la profundidad de empotramiento; los anclajes alivian
los esfuerzos estructurales pero no modifican la trayectoria hidráulica del agua.

Con el diseño actual ($D = 17{,}62$~m) se cumple $FS_\text{sifon} = 3{,}38 > 3$
y el margen respecto al límite hidráulico es de tan solo
$17{,}62 - 15{,}65 = \mathbf{1{,}97}$~m, lo que confirma que el sifonamiento
es la condición más restrictiva para cualquier reducción futura del
empotramiento.

\subsubsection{Ventajas e inconvenientes de la pantalla}

\begin{description}
  \item[Ventajas:]
    \begin{itemize}
      \item No requiere excavación previa: la pantalla se construye desde
            la rasante existente, minimizando las afecciones en superficie.
      \item Excelente impermeabilidad si las juntas entre paneles se ejecutan
            correctamente.
      \item Menor ocupación de suelo en planta respecto al muro ménsula
            (sin talón).
    \end{itemize}
  \item[Inconvenientes:]
    \begin{itemize}
      \item Requiere maquinaria especializada de gran tonelaje (hidrofresa), 
            no siempre disponible en obra de tamaño medio.
      \item Los lodos bentoníticos suponen un residuo de gestión controlada,
            incompatible con la \textbf{zona medioambientalmente protegida}
            en caso de vertido accidental.
      \item Necesidad de anclajes activos que invaden el trasdós; en una EDAR
            con instalaciones subterráneas esto puede ser inviable.
    \end{itemize}
\end{description}

% -----------------------------------------------------------
\subsection{Estimación de costes comparativa}
\label{subsec:costes_comparativa}
% -----------------------------------------------------------

Se realiza una estimación de orden de magnitud para los dos conceptos que
dominan el coste de cada solución: hormigón, acero y excavación. Los precios
unitarios aplicados son orientativos (HA-30: 105~€/m³; acero B~500~SD: 1{,}40~€/kg;
excavación convencional: 3~€/m³; excavación con lodos bentoníticos: 30~€/m³;
pilotes: 50~€/m³ excavación + 105~€/m³ hormigón + 1{,}40~€/kg acero).

\begin{notabox}[Hipótesis sobre la cuantía de acero]
  Se adopta una cuantía de \textbf{120~kg/m³} para la pantalla y de
  \textbf{90~kg/m³} para el muro en ménsula. La pantalla requiere mayor
  densidad de armado porque trabaja como elemento en voladizo de gran longitud
  ($L = 24{,}92$~m) sometido a momentos flectores muy superiores a los del fuste
  del muro en ménsula ($H = 7{,}30$~m empotrado en la zapata).
\end{notabox}

\subsubsection{Opción B — Muro pantalla}

Volumen total de hormigón (sección $0{,}75$~m $\times$ longitud $24{,}92$~m $\times$ 150~ml):
\[
  V_\text{pantalla} = 0{,}75 \times 24{,}92 \times 150 = \resultado{2\,803{,}5~\text{m}^3}
\]

\begin{align*}
  C_\text{HA-30}     &= 2\,803{,}5 \times 105 = 294\,367~\text{€} \\
  C_\text{acero}     &= 2\,803{,}5 \times 120~\text{kg/m}^3 \times 1{,}40
                     = 470\,988~\text{€} \\
  C_\text{exc.lodos} &= 2\,803{,}5 \times 30 = 84\,105~\text{€}
\end{align*}

\begin{table}[H]
  \centering
  \caption{Presupuesto estimado — Muro pantalla (150~ml)}
  \label{tab:presup_pantalla}
  \begin{tabular}{@{}lc@{}}
    \toprule
    \textbf{Concepto} & \textbf{Importe} \\
    \midrule
    Hormigón HA-30 ($2\,803{,}5$~m³ $\times$ 105~€/m³)   & 294\,367~€ \\
    Acero B~500~SD (120~kg/m³ $\times$ 1{,}40~€/kg)       & 470\,988~€ \\
    Excavación con lodos ($\times$ 30~€/m³)                &  84\,105~€ \\
    \midrule
    \textbf{TOTAL PANTALLA}                                & \textbf{849\,460~€} \\
    \bottomrule
  \end{tabular}
\end{table}

\subsubsection{Opción A — Muro en ménsula + pilotes}

Volumen de hormigón del muro (sección trapezoidal media $\approx e_2 \times H/2 \times B_z$):

\begin{align*}
  V_\text{muro} &= \left(e_1 + \frac{e_2-e_1}{2}\right) \cdot H \cdot L
                 + B \cdot h_z \cdot L \\
               &= \left(0{,}60 + \frac{0{,}15}{2}\right) \times 7{,}30 \times 150
                 + 5{,}60 \times 0{,}80 \times 150 \\
               &= 736{,}88 + 672{,}00 \approx \resultado{1\,411{,}1~\text{m}^3}
\end{align*}

\begin{align*}
  C_\text{HA-30, muro}    &= 1\,411{,}1 \times 105 = 148\,168~\text{€} \\
  C_\text{acero, muro}    &= 1\,411{,}1 \times 90~\text{kg/m}^3 \times 1{,}40
                          = 177\,799~\text{€} \\
  C_\text{exc. muro}      &= 1\,411{,}1 \times 3 = 4\,233~\text{€}
\end{align*}

Pilotes CPI ($D = 0{,}45$~m, $L = 0{,}5$~m), 1~pilote/ml, total 150~pilotes:
\[
  V_\text{pilote} = \frac{\pi \times 0{,}45^2}{4} \times 0{,}5
                  = \resultado{0{,}0795~\text{m}^3/\text{pilote}}
  \qquad
  V_\text{total} = 0{,}0795 \times 150 = \resultado{11{,}93~\text{m}^3}
\]

\begin{align*}
  C_\text{exc. pilotes}   &= 11{,}93 \times 50   =    596~\text{€} \\
  C_\text{HA-30, pilotes} &= 11{,}93 \times 105  =  1\,253~\text{€} \\
  C_\text{acero, pilotes} &= 11{,}93 \times 90 \times 1{,}40 = 1\,501~\text{€}
\end{align*}

\begin{table}[H]
  \centering
  \caption{Presupuesto estimado — Muro en ménsula + pilotes (150~ml)}
  \label{tab:presup_muro}
  \begin{tabular}{@{}lc@{}}
    \toprule
    \textbf{Concepto} & \textbf{Importe} \\
    \midrule
    Hormigón HA-30 muro ($1\,411{,}1$~m³ $\times$ 105~€/m³)  & 148\,168~€ \\
    Acero muro (90~kg/m³ $\times$ 1{,}40~€/kg)                & 177\,799~€ \\
    Excavación muro ($\times$ 3~€/m³)                          &   4\,233~€ \\
    Excavación pilotes ($\times$ 50~€/m³)                      &     596~€ \\
    Hormigón pilotes HA-30 ($\times$ 105~€/m³)                 &   1\,253~€ \\
    Acero pilotes (90~kg/m³ $\times$ 1{,}40~€/kg)              &   1\,501~€ \\
    \midrule
    \textbf{TOTAL MURO + PILOTES}                              & \textbf{333\,550~€} \\
    \bottomrule
  \end{tabular}
\end{table}

\begin{resultadobox}[Comparativa de costes]
  \begin{center}
  \begin{tabular}{lc}
    \toprule
    \textbf{Solución} & \textbf{Coste estimado} \\
    \midrule
    Pantalla anclada            & 849\,460~€ \\
    Muro en ménsula + pilotes   & 333\,550~€ \\
    \midrule
    \textbf{Ahorro con muro}    & \textbf{515\,910~€ ($\approx$61\,\%)} \\
    \bottomrule
  \end{tabular}
  \end{center}
\end{resultadobox}

% -----------------------------------------------------------
\subsection{Matriz de decisión comparativa}
\label{subsec:matriz_decision}
% -----------------------------------------------------------

\begin{table}[H]
  \centering
  \caption{Comparativa técnico-económica de las dos opciones}
  \label{tab:comparativa_opciones}
  \resizebox{\textwidth}{!}{%
  \begin{tabular}{@{}lcc@{}}
    \toprule
    \textbf{Criterio} & \textbf{Muro ménsula (HA)} &
    \textbf{Muro pantalla (anclada)} \\
    \midrule
    Altura de contención $H$ & 7{,}30~m \checkmark & 7{,}30~m \checkmark \\
    Factor seg.\ vuelco & 2{,}73 \checkmark & N/A (no aplica) \\
    Factor seg.\ desliz. & 1{,}88 \checkmark & N/A (no aplica) \\
    Solución de cimentación & Zapata + pilotes ($D=0{,}45$~m, $L=0{,}5$~m) & Pantalla (empotram.) \\
    Profundidad total de trabajo & 8{,}10~m (zapata) & 24{,}92~m (pantalla) \\
    Hormigón (m³) & $1\,411$~m³ (90~kg/m³ acero) & $2\,804$~m³ (120~kg/m³ acero) \\
    Coste estimado (150~ml) & \cellcolor{green!10}333\,550~€ & \cellcolor{red!10}849\,460~€ \\
    Ahorro & \multicolumn{2}{c}{515\,910~€ ($\approx$ 61\,\% menos con el muro)} \\
    Plazo estimado & $\sim$12 semanas & $\sim$18--20 semanas \\
    N.F.\ a $-2{,}50$~m & Drenaje de trasdós & Lodos + anclajes \\
    Sulfatos (3\,500~mg/kg) & HA-30/SR \checkmark & HA-30/SR + lodos SR \\
    Zona medioambiental prot. & Sin riesgo de vertido & Riesgo lodos \\
    Maquinaria & Convencional & Especializada (cuchara/hidrofresa) \\
    Mantenimiento & Sencillo (dren inspecc.) & Anclajes (retesado 5~años) \\
    Recomendación & \cellcolor{green!20}\textbf{RECOMENDADA} &
    \cellcolor{red!10}Descartada \\
    \bottomrule
  \end{tabular}}
\end{table}

% -----------------------------------------------------------
\subsection{Justificación de la solución adoptada}
\label{subsec:justificacion}
% -----------------------------------------------------------

\begin{supuestobox}[Solución adoptada: Muro de contención en ménsula de HA]
  Del análisis precedente se concluye que el \textbf{muro de contención
  en ménsula de hormigón armado} es la solución técnica y
  económicamente óptima para las condiciones de este proyecto. Los
  argumentos determinantes son:

  \begin{enumerate}
    \item \textbf{Estabilidad estructural verificada.} Las comprobaciones de
          vuelco ($FS = 2{,}73$) y deslizamiento ($FS = 1{,}88$) superan
          los límites reglamentarios. La capacidad portante del terreno de
          apoyo es suficiente. Los asientos estimados ($S_t = 7{,}51$~cm)
          superan el límite de $5$~cm, por lo que se adopta cimentación
          con \textbf{pilotes} de $D = 0{,}45$~m y $L = 0{,}5$~m
          (sección~\ref{subsec:pilotes_muro}), que reducen el asiento
          de grupo a $S_g = 0{,}66$~cm, despreciable.

    \item \textbf{Compatibilidad con el nivel freático.} El N.F.\ a
          $-2{,}50$~m se gestiona con un sistema de drenaje del trasdós
          (tubo dren $\varnothing\,100$ + geotextil + grava). 

    \item \textbf{Compatibilidad con la agresividad por sulfatos.}
          La especificación HA-30/SR resuelve la clase de exposición QA
          sin complejidades adicionales.

    \item \textbf{Economía y plazo.} El coste estimado del muro más
          pilotes asciende a $\approx 334$~kEUR frente a $\approx 849$~kEUR
          de la pantalla, lo que supone un ahorro de $\approx 516$~kEUR
          ($\sim$61\,\%).

    \item \textbf{Sencillez constructiva y mantenimiento.} El muro ménsula
          no requiere maquinaria especializada ni operaciones de retesado
          de anclajes a lo largo de la vida útil.
  \end{enumerate}

  La pantalla anclada sería la alternativa a considerar únicamente si
  existieran \textbf{restricciones de espacio en planta} que impidieran
  el desarrollo del talón de la zapata ($B_\text{talón} = 4{,}05$~m),
  lo cual no ocurre en la presente EDAR.
\end{supuestobox}
